\begin{abstract}
Numbers in the text of a scientific paper are usually typed by hand from analysis output, and
they drift from it when the analysis is rerun or the text is edited. We describe \pv{}, a tool
that generates every reported value from the result tables an analysis writes, so the manuscript
cites a table cell (\texttt{\textbackslash val\{table/key/column\}}) instead of a typed number.
For papers that already contain typed numbers, a language model proposes which table cell each
number reports and the program accepts the proposal only if the value in that cell rounds to the
written text; the model is never asked whether a number is correct. We tested the approach on
the Results section of one manuscript, a docking benchmark with 307 decimal numbers, by
corrupting each number in its last digit. Matching numbers to tables by value alone caught 1\% of
the corruptions and matching by surrounding words caught 20\%. With Claude Sonnet proposing
cells, all 134 anchored corruptions were caught and no correct number was flagged. Open models
run locally were not reliable when choosing from the whole table catalogue (Qwen2.5-32B proposed
139 wrong cells for 68 right ones) but agreed with Claude on 98\% of numbers when choosing only
among cells whose value already matched. Across two migrations, 312 of the 369 decimal numbers in
that Results section now come from data, and the compiled text of the paper did not change.
On a gold standard of 311 numbers from the migrated text, the best of eight local models
(Qwen3-32B) found the right cell for 48.9\% in the blind setting and 88.1\% when choosing among
matching cells, where wrong choices cannot be caught by value and still need review. \pv{} is open source under the MIT licence.
\end{abstract}
